\begin{proposition}[时间依赖张量的拉回微分]
设 $\Phi_{t,s}$ 为光滑时间依赖向量场 $V_t$ 的两参数局部流，$T_t$ 为光滑张量族。则\[\frac{\dd}{\dd t}(\Phi_{t,s}^*T_t)=\Phi_{t,s}^*(\partial_tT_t+\mathcal L_{V_t}T_t).\]
\end{proposition}

\begin{proof}
利用两参数流的组合律将 $\Phi_{t+\varepsilon,s}^*$ 分解为 $\Phi_{t,s}^*\circ\Phi_{t+\varepsilon,t}^*$，对短时流 $\Phi_{t+\varepsilon,t}^*$ 用 Lie 导数的一阶展开。

组合律分解。固定 $t$，取小量 $\varepsilon$。两参数流的组合律给出
\begin{equation*}
 \Phi_{t+\varepsilon,s}
 =\Phi_{t+\varepsilon,t}\circ\Phi_{t,s}.
\end{equation*}
因此拉回分解为
\begin{equation*}
 \Phi_{t+\varepsilon,s}^*T_{t+\varepsilon}
 =\Phi_{t,s}^*
 \left(\Phi_{t+\varepsilon,t}^*T_{t+\varepsilon}\right).
\end{equation*}

短时流展开。在时间 $t$ 附近，短时流 $\Phi_{t+\varepsilon,t}$ 由向量场 $V_t$ 生成，其拉回的一阶展开为
\begin{equation*}
 \Phi_{t+\varepsilon,t}^*
 =I+\varepsilon\,\mathcal L_{V_t}+o(\varepsilon),
\end{equation*}
而张量场的显式时间变化为
\begin{equation*}
 T_{t+\varepsilon}
 =T_t+\varepsilon\,\partial_tT_t+o(\varepsilon).
\end{equation*}

相乘并取差商。
\begin{align*}
 \Phi_{t+\varepsilon,t}^*T_{t+\varepsilon}
 &=\bigl(I+\varepsilon\,\mathcal L_{V_t}+o(\varepsilon)\bigr)
  \bigl(T_t+\varepsilon\,\partial_tT_t+o(\varepsilon)\bigr)\\
 &=T_t+\varepsilon
  \bigl(\partial_tT_t+\mathcal L_{V_t}T_t\bigr)+o(\varepsilon).
\end{align*}
因此
\begin{align*}
 \frac{\dd}{\dd t}\bigl(\Phi_{t,s}^*T_t\bigr)
 &=\lim_{\varepsilon\to0}\frac{\Phi_{t+\varepsilon,s}^*T_{t+\varepsilon}-\Phi_{t,s}^*T_t}{\varepsilon}\\
 &=\Phi_{t,s}^*
 \bigl(\partial_tT_t+\mathcal L_{V_t}T_t\bigr),
\end{align*}
即\eqnrefs{eq:math-dg-transport-identity}。
\end{proof}

\begin{proposition}[Levi--Civita 联络的度量变分]
设 $g_t$ 为光滑度量族，$h=\dot g$、$\mathcal A=\partial_t\nabla^{g_t}$。则 $\mathcal A$ 为张量，且\[\begin{aligned}2g(\mathcal A(X,Y),Z)={}&(\nabla_Xh)(Y,Z)+(\nabla_Yh)(X,Z)\\&-(\nabla_Zh)(X,Y).\end{aligned}\]
\end{proposition}

\begin{proof}
在每个固定 $t$ 上使用 Koszul 公式。向量场和 Lie 括号不随 $t$ 变化，所以对时间求导时只有度量与联络贡献。对
\[
2g(\nabla_XY,Z)
=Xg(Y,Z)+Yg(Z,X)-Zg(X,Y)
+g([X,Y],Z)-g([Y,Z],X)+g([Z,X],Y)
\]
求导，左边为
\[
2\dot g(\nabla_XY,Z)+2g(\mathcal A(X,Y),Z).
\]
右边把每个 $g$ 替换为 $\dot g$。再用当前 Levi--Civita 联络把普通方向导数拆成 $\nabla \dot g$ 与对向量槽位的联络作用；所有含 $\dot g(\nabla_XY,Z)$ 一类的项与左边对应项相消，得到
\begin{equation}\label{eq:math-dg-connection-variation}
{
\begin{aligned}
2g(\mathcal A(X,Y),Z)
={}&(\nabla_X\dot g)(Y,Z)
+(\nabla_Y\dot g)(X,Z)\\
&-(\nabla_Z\dot g)(X,Y).
\end{aligned}}
\end{equation}
坐标中即
\begin{equation}\label{eq:math-dg-christoffel-variation}
{
\partial_t\Gamma^k{}_{ij}
=\frac12g^{k\ell}
\left(
\nabla_i \dot g_{j\ell}
+\nabla_j \dot g_{i\ell}
-\nabla_\ell \dot g_{ij}
\right).}
\end{equation}

逐项消去的细节。对右边第一项 $Xg(Y,Z)$ 求导，先用 $\dot g$ 替换 $g$ 得 $\dot g(Y,Z)$ 的方向导数 $X\dot g(Y,Z)$，再用 Levi--Civita 联络展开
\begin{equation*}
X\dot g(Y,Z)
=(\nabla_X\dot g)(Y,Z)
+\dot g(\nabla_XY,Z)+\dot g(Y,\nabla_XZ).
\end{equation*}
对其他两项作类似展开。含 $\dot g(\nabla_XY,Z)$ 的项与左边 $2\dot g(\nabla_XY,Z)$ 部分相消：右边三项给出 $+1+1+0=2$ 倍 $\dot g(\nabla_XY,Z)$（前两项含 $\nabla_XY$，第三项含 $\nabla_ZX$ 不贡献），与左边相消。剩余项整理即得协变导数形式。

张量性的直接验证。要确认 $\mathcal A(X,Y)$ 是张量，需验证它对 $X,Y$ 都是函数线性的。对 $f\in C^\infty(M)$，
\begin{equation*}
\mathcal A(fX,Y)=\partial_t(\nabla_{fX}Y)=\partial_t(f\nabla_XY)=f\,\partial_t(\nabla_XY)=f\,\mathcal A(X,Y),
\end{equation*}
因 $f$ 与 $t$ 无关。对第二槽须显式保留联络的 Leibniz 项：
\[
\mathcal A(X,fY)=\partial_t\bigl(X(f)Y+f\nabla_XY\bigr)
=f\mathcal A(X,Y),
\]
因为 $X(f)Y$ 时间无关。这印证了"两个联络之差是张量"的一般事实。

Ricci 流 $\partial_tg=-2\operatorname{Ric}$ 给出 $\dot g=-2\operatorname{Ric}$，代入\eqnrefs{eq:math-dg-christoffel-variation} 得
\begin{equation*}
\partial_t\Gamma^k{}_{ij}
=-g^{k\ell}\left(
\nabla_i\operatorname{Ric}_{j\ell}
+\nabla_j\operatorname{Ric}_{i\ell}
-\nabla_\ell\operatorname{Ric}_{ij}
\right).
\end{equation*}

Hodge 星同时依赖 $g^{-1}$ 与体积形式 $\operatorname{vol}_g$，故其变分须合并\eqnrefs{eq:math-dg-inverse-metric-variation}与\eqnrefs{eq:math-dg-volume-variation}，下框给出完整公式。
\end{proof}

\begin{proposition}[曲率与标量曲率的变分]
对 $h=\dot g$、$\mathcal A=\dot\nabla$，\[\dot R(X,Y)Z=(\nabla_X\mathcal A)(Y,Z)-(\nabla_Y\mathcal A)(X,Z),\]\[\dot{\mathcal R}=-\langle h,\operatorname{Ric}\rangle_g+\nabla^i\nabla^jh_{ij}-\Delta_{\mathrm{LB}}\operatorname{tr}_gh.\]
\end{proposition}

\begin{proof}
从曲率定义
\[
R(X,Y)Z
=\nabla_X\nabla_YZ-\nabla_Y\nabla_XZ-\nabla_{[X,Y]}Z
\]
对时间求导。取 $X,Y,Z$ 为时间无关向量场，且 $\partial_t\nabla=\mathcal A$，得到
\[
\begin{aligned}
\partial_tR(X,Y)Z
={}&\mathcal A(X,\nabla_YZ)+\nabla_X\mathcal A(Y,Z)\\
&-\mathcal A(Y,\nabla_XZ)-\nabla_Y\mathcal A(X,Z)
-\mathcal A([X,Y],Z).
\end{aligned}
\]
另一方面，
\[
(\nabla_X\mathcal A)(Y,Z)
=\nabla_X(\mathcal A(Y,Z))
-\mathcal A(\nabla_XY,Z)
-\mathcal A(Y,\nabla_XZ).
\]
对 $X,Y$ 两项相减，并使用 Levi--Civita 联络无挠，
\[
\nabla_XY-\nabla_YX=[X,Y],
\]
便逐项恢复前一展开，从而得到\eqnrefs{eq:math-dg-curvature-variation-abstract}。在坐标基上取迹即得\eqnrefs{eq:math-dg-ricci-variation}。

最后对 $\mathcal R=g^{jk}\operatorname{Ric}_{jk}$ 求导：
\[
\partial_t\mathcal R
=(\partial_tg^{jk})\operatorname{Ric}_{jk}
+g^{jk}\partial_t\operatorname{Ric}_{jk}.
\]
第一项由\eqnrefs{eq:math-dg-inverse-metric-variation}给出
$-\langle\dot g,\operatorname{Ric}\rangle_g$。第二项把\eqnrefs{eq:math-dg-christoffel-variation}代入\eqnrefs{eq:math-dg-ricci-variation}，利用 $\nabla g=0$ 收缩指标，得到
\[
g^{jk}\partial_t\operatorname{Ric}_{jk}
=\nabla^i\nabla^j\dot g_{ij}
-\Delta_{\mathrm{LB}}(\operatorname{tr}_g\dot g),
\]
两项相加即为\eqnrefs{eq:math-dg-scalar-curvature-variation}。

标量曲率变分公式\eqnrefs{eq:math-dg-scalar-curvature-variation} 是几何流理论的基础：代入 Ricci 流 $\dot g=-2\operatorname{Ric}$ 得 $\partial_t\mathcal R=\Delta\mathcal R+2|\operatorname{Ric}|^2$，即标量曲率沿 Ricci 流 的反应--扩散方程。

共形变分 $\dot g=2u\,g$ 下，$\operatorname{tr}_g\dot g=2nu$、$\nabla^i\nabla^j\dot g_{ij}=2\Delta u$，代入得 $\partial_t\mathcal R=-2(n-1)\Delta_{\mathrm{LB}}u-2u\mathcal R$；$n=2$ 时为 $\partial_t\mathcal R=-2\Delta_{\mathrm{LB}}u-2u\mathcal R$，即 Gauss 曲率的共形演化。
\end{proof}

\begin{theorem}[固定体积的总标量曲率临界点]
设 $M^n$ 紧致、无边界，$n\ge3$。总标量曲率 $\mathcal S[g]=\int_M\mathcal R_g\operatorname{vol}_g$ 在固定总体积的度量空间上的临界点恰为 Einstein 度量 $\operatorname{Ric}=cg$。二维时该泛函由 Gauss--Bonnet 固定，不能由此推出常曲率。
\end{theorem}

\begin{proof}
对体积泛函
\[
\mathcal V[g]=\int_M\operatorname{vol}_g
\]
有
\[
\delta\mathcal V_g[h]
=\frac12\int_M\operatorname{tr}_g(h)\,\operatorname{vol}_g
=\frac12\int_M\langle g,h\rangle_g\operatorname{vol}_g.
\]
临界条件
\[
\delta(\mathcal S+2\lambda\mathcal V)[h]=0
\quad\forall h
\]
给出
\[
-\operatorname{Ric}+\frac12\mathcal R g+\lambda g=0.
\]
在 $n$ 维中取迹：
\[
-\mathcal R+\frac n2\mathcal R+n\lambda=0,
\]
$n\ge3$ 时可解出 $\mathcal R=-2n\lambda/(n-2)$，为常数。二维时这个取迹式只有 $\lambda=0$，并不能推出 $\mathcal R$ 恒定。代回原方程，得到
\[
\operatorname{Ric}=\left(\frac12\mathcal R+\lambda\right)g=:cg.
\]
反过来任何 Einstein 度量 满足带适当乘子的 Euler--Lagrange 方程。

体积变分的显式推导。度量变分 $g_t=g+th$ 下体积形式的变化为
\begin{equation*}
\operatorname{vol}_{g_t}
=\sqrt{\frac{\det(g+th)}{\det g}}\,\operatorname{vol}_g
=\sqrt{\det(I+tg^{-1}h)}\,\operatorname{vol}_g.
\end{equation*}
对 $t$ 在 $t=0$ 求导，用 $\det(I+A)$ 的一阶展开 $\det(I+A)=1+\operatorname{tr}A+O(\|A\|^2)$，
\begin{equation*}
\left.\frac{\dd}{\dd t}\operatorname{vol}_{g_t}\right|_{t=0}
=\frac12\operatorname{tr}(g^{-1}h)\,\operatorname{vol}_g
=\frac12\operatorname{tr}_g(h)\,\operatorname{vol}_g.
\end{equation*}
积分即得体积泛函的变分公式。$\operatorname{tr}_g(h)=g^{ij}h_{ij}$ 是关于 $g$ 的迹，与坐标无关。

常数 $c$ 的符号给出 Ricci 曲率的分类：$c>0$ 为正 Ricci（例如球面），$c=0$ 为 Ricci 平坦，$c<0$ 为负 Ricci（例如双曲空间）。Einstein 条件在维数大于三时不意味着所有截面曲率相同，不能据 $c$ 的符号直接分类全部截面曲率。Einstein 度量还是 Ricci 流 的缩放固定点（$g(t)=(1-2ct)g$），并在广义相对论中对应带宇宙学常数的真空 Einstein 方程。
\end{proof}

\begin{proposition}[Hodge 星的度量变分]
设 $g_t$ 为定向流形上的光滑度量族，$\mathcal V=g^{-1}\dot g$，$\alpha$ 为时间无关 $k$-形式。则\[\partial_t(*\alpha)=*\left[\tfrac12\operatorname{tr}_g\dot g\,\alpha-\mathcal V^{[k]}\alpha\right].\]特别地，$\dot g=2ug$ 时 $\partial_t(*\alpha)=(n-2k)u*\alpha$。
\end{proposition}

\begin{proof}
Hodge 星由度量通过内积和体积形式隐式定义。对度量变分，分别追踪内积变化和体积变化，再用 Hodge 星定义反读。

Hodge 星的隐式定义。对任意时间无关 $k$-形式 $\alpha,\beta$，Hodge 星由
\begin{equation*}
 \alpha\wedge *\beta
 =\langle\alpha,\beta\rangle_g\,\operatorname{vol}_g
\end{equation*}
唯一确定。两边对 $t$ 求导，$\alpha$ 时间无关故只需变分 $*\beta$、$\langle\cdot,\cdot\rangle_g$ 和 $\operatorname{vol}_g$。

内积变分。逆度量变化\eqnrefs{eq:math-dg-inverse-metric-variation}意味着 $k$-形式点态内积满足
\begin{equation*}
 \partial_t\langle\alpha,\beta\rangle_g
 =-\langle \mathcal V^{[k]}\alpha,\beta\rangle_g,
\end{equation*}
其中 $\mathcal V^{[k]}$ 是度量变分算子 $\mathcal V=g^{-1}\dot g$ 在 $k$-形式上的自然提升。

体积变分。由\eqnrefs{eq:math-dg-volume-variation}，
\begin{equation*}
 \partial_t\operatorname{vol}_g
 =\frac{1}{2}\operatorname{tr}_g(\dot g)\,\operatorname{vol}_g.
\end{equation*}

合并。对 Hodge 星的隐式定义等式求导，
\begin{align*}
 \alpha\wedge\partial_t(*\beta)
 &=\left[\partial_t\langle\alpha,\beta\rangle_g\right]\operatorname{vol}_g
  +\langle\alpha,\beta\rangle_g\,\partial_t\operatorname{vol}_g\\
 &=\left[
  -\langle \mathcal V^{[k]}\alpha,\beta\rangle_g
  +\frac{1}{2}\operatorname{tr}_g(\dot g)
  \langle\alpha,\beta\rangle_g
 \right]\operatorname{vol}_g.
\end{align*}

自伴性。因为 $\mathcal V$ 自伴（$g^{-1}\dot g$ 关于 $g$ 对称），诱导的 $\mathcal V^{[k]}$ 在形式内积下也自伴：
\begin{equation*}
 \langle \mathcal V^{[k]}\alpha,\beta\rangle_g
 =\langle\alpha,\mathcal V^{[k]}\beta\rangle_g.
\end{equation*}

反读 Hodge 星。将自伴性代入上面合并得到的等式，
\begin{equation*}
 \alpha\wedge\partial_t(*\beta)
 =\left\langle
  \frac{1}{2}\operatorname{tr}_g(\dot g)\,\alpha
  -\mathcal V^{[k]}\alpha,\,
  \beta
 \right\rangle_g\operatorname{vol}_g
 =\alpha\wedge *\left[
  \frac{1}{2}\operatorname{tr}_g(\dot g)\,\beta
  -\mathcal V^{[k]}\beta
 \right],
\end{equation*}
由 $\alpha$ 任意性即得到\eqnrefs{eq:math-dg-hodge-star-variation}。

共形核对：$\dot g=2ug$ 时 $\mathcal V=2uI$，
其在 $k$-形式上的提升为 $\mathcal V^{[k]}=2kuI$，故
\[
\dot *\alpha=(n-2k)u*\alpha.
\]
这与 $g\mapsto \ee^{2f}g$ 时 $*\mapsto \ee^{(n-2k)f}*$ 相容；
中间次数 $k=n/2$ 的 Hodge 星共形不变。
\end{proof}

\begin{proposition}[Ricci 孤子产生自相似流]
若 $\operatorname{Ric}_g+\frac12\mathcal L_Xg=\lambda g$，令 $c(t)=1-2\lambda t>0$，$\phi_t$ 由 $X/c(t)$ 生成。则在流存在的时间区间，$g(t)=c(t)\phi_t^*g$ 精确满足 Ricci 流方程。
\end{proposition}

\begin{proof}
设 $(g,X,\lambda)$ 满足\eqnrefs{eq:math-dg-ricci-soliton}，令
\[
c(t)=1-2\lambda t.
\]
在 $c(t)>0$ 的区间上，令 $\phi_t$ 由时间依赖向量场
\[
Y_t=\frac{1}{c(t)}X
\]
生成，并定义
\begin{equation}
g(t)=c(t)\phi_t^*g.
\label{eq:math-dg-soliton-selfsimilar-flow}
\end{equation}
则
\begin{align*}
\partial_tg(t)
&=c'(t)\phi_t^*g
+c(t)\phi_t^*(\mathcal L_{Y_t}g)\\
&=\phi_t^*(-2\lambda g+\mathcal L_Xg)\\
&=-2\phi_t^*\operatorname{Ric}_g.
\end{align*}
常数缩放不改变 $(0,2)$ 型 Ricci 张量，而 Ricci 张量 对微分同胚自然，因此
\[
\operatorname{Ric}_{g(t)}=\phi_t^*\operatorname{Ric}_g.
\]
所以
\[
\partial_tg(t)=-2\operatorname{Ric}_{g(t)}.
\]
这证明 孤子 不是"近似自相似"，而是 Ricci 流 在缩放与微分同胚作用下的精确固定点。
\end{proof}
